% §I Introduction (~1.0 page, Fig. 1). REWRITE 08-11 (사용자 결정): prescription-led arc —
% problem → cause (bypass) → recipe (grounded bottleneck + decorrelated demos) → thesis.
% Thesis = crossover: "indistinguishable in loss, similar at demonstrated forces, categorically
% different exactly where the task is decided." Quintuple recast as ablation of the recipe.
% Robot evidence in force language (contact-force distributions, censoring); success counts
% deferred to §VI and subordinated there. Nugget sentence PROMOTED here from old §VI closer —
% §VI rewrite must not repeat it verbatim.
% Forbidden narratives respected: no "naive is force-blind" (state-transplant 105%); scope
% limited to correlated-BC + existing shortcut (토론 1/5). All numbers EVIDENCE §3, 0729-pure.
% Terminology 08-11 (사용자): unconditioned policy = "naive policy" (defined at first mention
% with "full wrench in its state"), replacing "baseline" PAPER-WIDE — rename §III–§VII prose
% and table labels ("naive VLA (full wrench)") as each section is rewritten.
\section{Introduction}
\label{sec:intro}

Behavior-cloned vision-language-action (VLA) policies are rapidly acquiring physical condition
signals. Wrist force/torque, contact events, and task-state bits now routinely appear in the
observation vector, and a growing line of work fuses force into the policy backbone for
contact-rich manipulation~\cite{forcevla,phaforce,forceflow,filic}. All of this work shares a
validation recipe---add the signal, retrain, and validate by task success (occasionally
alongside rollout force statistics~\cite{phaforce,forceflow})---and none of it measures
whether the trained policy \emph{causally uses} the signal when it matters: whether, holding
everything else fixed, changing the force input changes the action. Imitation learning offers no
such guarantee. The objective asks the policy to reproduce \emph{what} the expert did, not
\emph{why}: a clone can imitate the behavior convincingly while being driven by entirely
different cues than the expert was. Imitation pins down the trajectory, not the
mechanism---equally faithful clones can act from opposite causes, and the loss cannot tell
them apart~\cite{dehaan2019causal,geirhos2020shortcut}.

The consequence is easiest to see where force is the only signal that can carry the task. We
study real-robot suction case picking with a single head-mounted camera and no wrist camera: the
arm must stop descending the moment the suction cup meets the case, and stack heights differing
by a few millimeters---the entire success margin---are visually near-indistinguishable at
head-camera range. The wrench is in the state the policy receives, and every demonstration
exhibits the stop. Yet a naively fine-tuned VLA---full wrench in its state, hereafter the
\emph{naive policy}---presses through contact in every evaluation trial and never initiates a
stop of its own: each run ends only when an external safety layer, separate from the policy,
trips at its $15$\,N force limit---set to protect the battery cells the cases carry
($15.4$--$18.7$\,N measured at abort)---the policy still commanding descent. The task turns on a single transition---stop when force
rises---and it is precisely this transition that the policy fails to bind to force.

Our fix changes how force is delivered, not what is delivered. A condition vector \chat{}
(contact, normal force, seal) is \emph{computed} from signals the naive policy already receives
and injected through FiLM modulation~\cite{perez2018film}; the raw wrench is masked from the
state so that the same signal does not enter twice, making \chat{} the \emph{only} pathway from
force to action. We call this pathway \emph{grounded} because its content is computed from
measured signals, and the policy trained behind it the \emph{conditioned policy} (``FiLM
(ours)'' in tables and figures). By
construction it receives no information the naive policy lacks, and because the pathway is
explicit and unique, forcing \chat{} at inference directly measures its \emph{causal
authority}---the fraction of commanded descent the condition can cancel.

The mask is load-bearing. Leave the raw wrench in the state and the conditioning pathway trains
causally dead at unchanged validation loss: training routes force through the already-functional
raw path and starves the dedicated one. We call this \emph{bypass} and exhibit it in
Sec.~\ref{sec:offline}; it is why a fix of this kind cannot be certified after the fact, and why
ours holds by construction. The demonstrations complete the recipe: each press-retreat episode
presses to a randomized force target and retreats before any seal event, so the force at which
the expert stops varies from episode to episode---in the data, depth, appearance, and timing no
longer predict the stop; only the force reading does. The same interface that fixes the policy
is the instrument that audits it (Sec.~\ref{sec:method}).

The resulting policy is not uniformly better---it is different exactly where the task is
decided. In a matched comparison---same data, same initialization, imitation error
indistinguishable---the naive policy \emph{does} respond to force at the demonstrated
magnitudes, as strongly as the conditioned policy or more so. But its response is a
\emph{template}: sized to the demonstrated forces, it fades or saturates as force rises past
them, never strengthening into a stop. The conditioned policy's response instead grows with
force, to the point of stopping and
backing away at forces beyond anything in the training data (Sec.~\ref{sec:offline}). The
mechanism carries to the real closed loop: while the naive policy drives contact force past the
$15$\,N safety limit on every attempt, the conditioned policy stops on its own in every
pick---across three stack heights, at a median contact force of $2.5$\,N---and live
counterfactual probes reproduce the offline crossover on the physical system
(Sec.~\ref{sec:robot}).

Our contribution is the pathway and three properties of it, each measured by intervention on
the trained policy (Sec.~\ref{sec:method-probes}):
\begin{itemize}
\item \textbf{Drop-in.} \chat{} is computed from signals the policy already receives and
  enters through an identity-initialized FiLM layer at the state token; the backbone is
  untouched, no information is added, imitation error is unchanged, and the same pathway
  attaches to a second backbone ($\pi_0$), where it again turns a policy that never stops into
  one that does (Sec.~\ref{sec:method},~\ref{sec:offline-pi0}).
\item \textbf{Retrofittable.} Added to an already-trained naive policy and continued for a
  few thousand steps on the same demonstrations, the pathway turns a saturating force template
  into a monotone brake; a matched ablation shows the mask and the grounding are both
  necessary---unmasked, the pathway trains causally dead (\emph{bypass}); ungrounded, it is the
  worst policy in the study (Sec.~\ref{sec:offline}).
\item \textbf{Monotone beyond the demonstrations.} Because \chat{} is a calibrated monotone
  scalar, braking grows with force past anything demonstrated, stopping and retreating where
  the naive policy's template saturates and presses through (Sec.~\ref{sec:offline}). The
  same holds on our robot: every conditioned pick stops under the $15$\,N abort limit (median
  contact force $2.5$\,N), including at a stack height absent from the demonstrations, and
  live counterfactuals reproduce the crossover (Sec.~\ref{sec:robot}).
\end{itemize}

\noindent In one sentence: \emph{a policy that stops when it feels what the demonstrations felt
and a policy that stops harder the harder it is pressed are indistinguishable in loss and
similar at the demonstrated forces---they differ exactly where the task is decided, and only
demonstrations in which force alone predicts the stop, behind a grounded, masked FiLM
pathway, yield the second.}
